% Generated draft: requires mathematical and source review.
\documentclass[11pt]{article}
\usepackage[margin=25mm]{geometry}
\usepackage{fontspec}
\setmainfont{Latin Modern Roman}
\usepackage{amsmath,amssymb,mathtools,unicode-math}
\setmathfont{Latin Modern Math}
\usepackage{enumitem,needspace,xurl,hyperref}
\setlength{\parindent}{0pt}
\setlength{\parskip}{5pt}
\setlength{\emergencystretch}{3em}
\newcommand{\N}{\mathbb{N}}
\newcommand{\Z}{\mathbb{Z}}
\newcommand{\Q}{\mathbb{Q}}
\newcommand{\R}{\mathbb{R}}
\newcommand{\C}{\mathbb{C}}
\newcommand{\F}{\mathbb{F}}
\newcommand{\A}{\mathbb{A}}
\newcommand{\PP}{\mathbb{P}}
\newcommand{\E}{\mathbb{E}}
\newcommand{\eps}{\varepsilon}
\newcommand{\dd}{\,\mathrm{d}}
\begin{document}
\section*{UnsolvedMath ID: 6400008}
\textbf{Source problem code:} catalogue entry
\subsection{Definitions and mathematical statement}
Let $M$ be a compact smooth seven-manifold admitting the source's holonomy-$G_2$ metrics. To make the proposed cohomology map well defined, one normally works with torsion-free $G_2$-structures: positive three-forms $\varphi$ inducing a metric $g_\varphi$ and satisfying
\[
 \mathrm d\varphi=0,\qquad
 \mathrm d(*_\varphi\varphi)=0.
\]
An orientation and positivity convention are understood in this formulation. Whether holonomy must be exactly $G_2$ or may be a subgroup of $G_2$ is not fixed in the supplied wording.

Let $\mathcal F$ be the intended space of such forms and let $\operatorname{Diff}_0(M)$ be diffeomorphisms isotopic to the identity. The structure-moduli interpretation is
\[
 \mathcal M=\mathcal F/\operatorname{Diff}_0(M),\qquad
 \pi:\mathcal M\to H^3(M;\R),\quad
 \pi([\varphi])=[\varphi]_{\mathrm{dR}}.
\]
The cohomology class is unchanged by diffeomorphisms isotopic to the identity. The source calls this a moduli space of metrics and speaks of the metric's defining form; a metric alone does not include all the needed form/orientation choices in every holonomy convention. The form-moduli versus metric-moduli distinction must therefore be resolved.

The questions are to describe this moduli space and, assuming the stated local-diffeomorphism property, to determine whether $\pi$ is globally a diffeomorphism onto its image. In the structure-moduli formulation the proposed conclusion is that $\pi$ is injective on the whole intended $\mathcal M$, so its local smooth inverses combine to a global smooth inverse on $\pi(\mathcal M)$. Equivalently, equal degree-three cohomology classes should force two admissible structures to differ by a diffeomorphism isotopic to the identity. The question does not ask surjectivity onto all of $H^3(M;\R)$.

The original wording does not restrict this injectivity to a single connected component. Componentwise results or a quotient by the full diffeomorphism group would be different statements.
\subsection{Short English statement}
Describe the global moduli of compact $G_2$ geometry and decide whether its degree-three period map is globally one-to-one, given that it is locally a diffeomorphism. The distinction between metrics and defining torsion-free forms is essential.

The supplied summary attributes local smoothness, harmonic-three-form tangent spaces, and the local period-map result to Joyce and Hitchin. It calls the general global question open, but also mentions an open-embedding claim on components and Calabi--Yau Torelli-type comparisons without exact hypotheses or references. Its final sentence is truncated. Those comments are not adopted as a general injectivity theorem. The partial-progress label and 2008 source are retained with current status marked for review.
\subsection{Sources}
\begin{itemize}
\item Local UnsolvedMath registry, permanent ID 6400008, source code catalogue entry, ``Global Moduli of G2 Metrics,'' proposed by Simon Donaldson, 2008. Dataset curation and metadata: Ulam AI and UnsolvedMath Contributors, CC BY 4.0.
\item Simon Donaldson, \emph{Some Problems in Differential Geometry and Topology} (2008), explicit question/program 8 in the author PDF, \url{https://wwwf.imperial.ac.uk/~skdona/NONLINEARITYPROBLEMS.PDF}; normalized pdf-text extraction, accessed 2026-07-29.
\item Supplied summary attributions: Joyce and Hitchin, for local $G_2$ moduli and period-map results; unspecified restricted-class and Calabi--Yau Torelli comparisons. No full bibliographic entries or exact global theorem are supplied.
\item Source-specific provenance marks current status and statement-reuse terms NEEDS\_REVIEW for the public author-hosted PDF, separately from the dataset licence.
\end{itemize}
\end{document}
